% !TEX program = xelatex
\documentclass[11pt,a4paper]{article}
% !TEX program = xelatex
\usepackage[a4paper,top=23mm,bottom=24mm,left=21mm,right=21mm,headheight=15pt]{geometry}
\usepackage{fontspec}
\usepackage{kotex}
\usepackage{graphicx}
\usepackage{booktabs}
\usepackage{array}
\usepackage{tabularx}
\usepackage{makecell}
\usepackage[inline]{enumitem}
\usepackage{setspace}
\usepackage{xcolor}
\usepackage{fancyhdr}
\usepackage{lastpage}
\usepackage{calc}
\usepackage{needspace}
\usepackage{hyperref}

\defaultfontfeatures{Ligatures=TeX}

% Korean exam style: body and headings are set in a true Myeongjo/Batang-style face.
% This avoids sans-serif fallback that can look like Gulim in PDF viewers.
\setmainfont{Linux Libertine O}
\setsansfont{Noto Sans}
\setmonofont{DejaVu Sans Mono}
\setmainhangulfont{NanumMyeongjo}
\setsanshangulfont{NanumBarunGothic}

% Linguistic forms and IPA.
\newfontfamily\lingfont{Charis SIL}[Scale=MatchLowercase]
\hypersetup{colorlinks=true, linkcolor=black, urlcolor=black, pdfborder={0 0 0}}
\definecolor{RuleGray}{HTML}{4F4F4F}
\definecolor{LightRuleGray}{HTML}{B7B7B7}

\setlength{\parindent}{0pt}
\setlength{\parskip}{0.36em}
\setlength{\tabcolsep}{0.62em}
\renewcommand{\arraystretch}{1.16}
\raggedbottom
\sloppy
\emergencystretch=2em

\pagestyle{fancy}
\fancyhf{}
\fancyhead[L]{\footnotesize\bfseries 2025-2학기 특별 세션 1}
\fancyhead[R]{\footnotesize\bfseries 미니 언어학 올림피아드}
\fancyfoot[C]{\footnotesize \thepage\,/\,\pageref*{LastPage}}
\renewcommand{\headrulewidth}{0.35pt}
\renewcommand{\footrulewidth}{0pt}

\fancypagestyle{plain}{%
  \fancyhf{}%
  \fancyfoot[C]{\footnotesize \thepage\,/\,\pageref*{LastPage}}%
  \renewcommand{\headrulewidth}{0pt}%
  \renewcommand{\footrulewidth}{0pt}%
}

\newcommand{\term}[1]{{\lingfont\bfseries #1}}
\newcommand{\ipa}[1]{{\lingfont [#1]}}
\newcommand{\blankko}[1]{(\textbf{#1})}
\newcommand{\subq}[1]{\par\smallskip\noindent\textbf{(#1)}\quad}
\newcommand{\credit}[1]{\nobreak\hspace{1em}\hfill{\bfseries - #1}\par}
\newcommand{\warnnote}{\par\smallskip\noindent{\bfseries 참고}\quad\ignorespaces}

\newcommand{\examtitle}{%
  \thispagestyle{plain}%
  \begin{center}
    {\bfseries\Large 2025-2학기 특별 세션 1\par}
    \vspace{0.25em}
    {\bfseries\huge 미니 언어학 올림피아드\par}
    \vspace{0.65em}
    {\large 서울대학교 언어학회 \quad\textbar\quad 2025년 10월 1일\par}
  \end{center}
  \vspace{0.7em}

  \noindent\rule{\textwidth}{0.4pt}
}

\newcommand{\answertitle}{%
  \thispagestyle{plain}%
  \begin{center}
    {\bfseries\Large 2025-2학기 특별 세션 1\par}
    \vspace{0.25em}
    {\bfseries\huge 미니 언어학 올림피아드: 답안지\par}
    \vspace{0.65em}
    {\large 서울대학교 언어학회 \quad\textbar\quad 2025년 10월 1일\par}
  \end{center}
  \vspace{0.7em}
  \noindent\rule{\textwidth}{0.8pt}
}

\newcommand{\problem}[1]{%
  \par\vspace{1.2em}\needspace{7\baselineskip}%
  \noindent{\bfseries\Large #1}\par
  \vspace{-0.1em}\noindent\textcolor{RuleGray}{\rule{\textwidth}{0.5pt}}\par\vspace{0.45em}
}

\newcommand{\answersection}[1]{%
  \par\vspace{1.1em}\needspace{6\baselineskip}%
  \noindent{\bfseries\Large #1}\par
  \vspace{-0.1em}\noindent\textcolor{RuleGray}{\rule{\textwidth}{0.5pt}}\par\vspace{0.45em}
}

\newcommand{\exampleblock}[4]{%
  \begin{minipage}[t]{0.47\textwidth}
  \raggedright\footnotesize
  \term{#1}\par\vspace{0.05em}
  {\lingfont #2}\par\vspace{0.05em}
  {\scriptsize #3}\par\vspace{0.05em}
  #4
  \end{minipage}%
}

\newcommand{\answerline}[2]{\textbf{#1}\quad #2\\}

% Deseret
\newfontfamily\deseretfont{Noto Sans Deseret}[Scale=1.22]

\newcommand{\des}[1]{{\deseretfont #1}}
\newcommand{\desword}[2]{%
  \makebox[2.7em][r]{(#1).}\,\des{#2}%
}


\begin{document}
\answertitle
\begin{spacing}{1.15}

\answersection{연습문제 정답}

\begin{center}
\begin{tabular}{@{}ll@{}}
\toprule
빈칸 & 정답 \\
\midrule
\blankko{ㄱ} & to squeeze between \\
\blankko{ㄴ} & \term{nəɣbər} \\
\blankko{ㄷ} & \term{mətəltəl} \\
\blankko{ㄹ} & to be tangled up \\
\blankko{ㅁ} & to be bothered by s.b. \\
\bottomrule
\end{tabular}
\end{center}

\vspace{0.5em}

\textbf{규칙.} 어간이 \term{yi-}로 시작하면 \term{yi-}를 삭제하고, 그렇지 않으면 그대로 둔 뒤 \term{mə-}를 붙인다. 어간에 입술소리 \term{m, b, f}가 포함되어 있으면 \term{mə-}가 \term{nə-}가 된다.

\answersection{문제 \#1 정답}
\begin{center}
\begin{tabular}{@{}ll@{}}
\toprule
로비아나어 & 한국어 \\
\midrule
\term{toa} & 살다 \\
\term{gani} & 먹다 \\
\term{kinera} & 노래 \\
\term{hiva} & 원하다 \\
\term{halehaleana} & 계단 \\
\term{zinamana} & 그의 언어 \\
\term{tinoaqu} & 나의 삶 \\
\term{kera} & 노래하다 \\
\term{hinivana} & 그의 소망 \\
\term{hale} & 오르다 \\
\term{habohabotuana} & 의자 \\
\term{tinoana} & 그의 삶 \\
\bottomrule
\end{tabular}
\end{center}

\subq{b} \term{habotu} = 앉다, \term{hinivaqu} = 나의 소망

\subq{c} 말하다 = \term{zama}, 음식 = \term{ginani}

\textbf{규칙.} \term{X-qu}는 `나의 X', \term{X-na}는 `그의 X'를 뜻한다. \term{C-in-V...}는 명사화소로 쓰이며, \term{동사-동사-ana}는 `X할 수 있는 물건'을 뜻한다.

\answersection{문제 \#2 정답}
\textbf{스페인어 : 무일락어 대응}

\begin{center}
\begin{tabular}{@{}ll@{}}
\toprule
스페인어 & 무일락어 \\
\midrule
\term{b, v} & \term{w} \\
\term{c} & \term{k} \\
\term{r, rr} & \term{ʀ} \\
\term{i, e} & \term{i} \\
\term{u, o} & \term{u} \\
\term{Cd} & \term{Ct} \\
\term{Vd} & \term{Vw} \\
\bottomrule
\end{tabular}
\end{center}

\subq{a} \term{sawadu} `토요일'이 예외이다. 예상되는 형태는 \term{sawawu}이다.

\subq{b}
\begin{center}
\begin{tabular}{@{}lll@{}}
\toprule
스페인어 & 의미 & 무일락어 \\
\midrule
\term{vela} & 양초 & \term{wila} \\
\term{condenado} & 유령 & \term{kuntinawu} \\
\term{barreta} & 작은 지레 & \term{waʀita} \\
\term{palabra} & 단어 & \term{palawʀa} \\
\term{corbata} & 넥타이 & \term{kuʀwata} \\
\term{basura} & 쓰레기 & \term{wasuʀa} \\
\bottomrule
\end{tabular}
\end{center}

\answersection{문제 \#3 정답}
\subq{a}
\begin{center}
\begin{tabular}{@{}llllllll@{}}
\toprule
\blankko{ㄱ} & 9 & \blankko{ㄴ} & 3 & \blankko{ㄷ} & 6 & \blankko{ㄹ} & 12 \\
\blankko{ㅁ} & 5 & \blankko{ㅂ} & 2 & \blankko{ㅅ} & 1 & \blankko{ㅇ} & 7 \\
\blankko{ㅈ} & 4 & \blankko{ㅊ} & 8 & \blankko{ㅋ} & 14 & \blankko{ㅌ} & 10 \\
\blankko{ㅍ} & 13 & \blankko{ㅎ} & 11 & & & & \\
\bottomrule
\end{tabular}
\end{center}

\subq{b} 가능한 표기:
\begin{center}
\large
\setlength{\tabcolsep}{0.6em}
\renewcommand{\arraystretch}{1.6}
\begin{tabular}{@{}ccccc@{}}
(15) \des{𐐻𐐴𐐿} &
(16) \des{𐑅𐐻𐑉𐐴𐐿} &
(17) \des{𐑋𐐲𐑇𐑉𐐭𐑋} &
(18) \des{𐐰𐑍𐑀𐑉𐐨} &
(19) \des{𐑇𐑉𐐮𐑍𐐿}
\end{tabular}
\end{center}

\subq{c} \blankko{ㄲ} = \ipa{hæpɪ}, \blankko{ㄸ} = \ipa{gʊd}

\answersection{문제 \#4 정답}
\begin{center}
\small
\begin{tabular}{@{}cl@{\qquad}cl@{}}
\toprule
빈칸 & 정답 & 빈칸 & 정답 \\
\midrule
(1) & \term{karenaumo} & (9) & \term{kivoraita kena:nka ta:} \\
(2) & \term{paʔnokoʔ keta be:naumo ta:} & (10) & \term{kena:nka kivoraita paʔnokoʔ ta:} \\
(3) & \term{paʔnokoʔ keta naruŋ ta:} & (11) & \term{nebolo aɰava} \\
(4) & \term{ela} 또는 \term{pa} & (12) & \term{iɰivanipiti} \\
(5) & \term{kena:nka kivoraita be:naumo ta:} & (13) & \term{poɰoanipiti} \\
(6) & \term{kivoraita paʔnokoʔ keta be:naumo ta:} & (14) & \term{kuavanipiti} \\
(7) & \term{kivoraita paʔnokoʔ keta kena:nka ta:} & (15) & \term{nebolo} \\
(8) & \term{kivoraita karenaumo ta:} & & \\
\bottomrule
\end{tabular}
\end{center}

3$\times$3 마방진은 다음 값 배열을 갖는다.
\[
\begin{array}{ccc}
4 & 9 & 2 \\
3 & 5 & 7 \\
8 & 1 & 6
\end{array}
\]
4$\times$4 마방진은 다음 값 배열을 갖는다.
\[
\begin{array}{cccc}
10 & 24 & 23 & 13 \\
21 & 15 & 16 & 18 \\
17 & 19 & 20 & 14 \\
22 & 12 & 11 & 25
\end{array}
\]

\answersection{문제 \#5 정답}
\begin{center}
\begin{tabular}{@{}ll@{}}
\toprule
빈칸 & 정답 \\
\midrule
\blankko{ㄱ} & 2 - 22 - 42 - 31 \\
\blankko{ㄴ} & 2 - 3 - 1 \\
\blankko{ㄷ} & 2 - 3 - 2 - 21 - 4 \\
\blankko{ㄹ} & 2 - 3 - 2 - 322 - 21 \\
\blankko{ㅁ} & 2 - 32 - 3 - 2 - 3 - 2 - 1 - 4 \\
\blankko{ㅂ} & 3 - 22 - 3 - 2 - 31 \\
\blankko{ㅅ} & 2 - 31 - 322 - 31 \\
\blankko{ㅇ} & 2 - 22 - 32 - 22 - 21 \\
\bottomrule
\end{tabular}
\end{center}

\textbf{규칙.} 1 < 2 < 3 < 4로, 숫자가 클수록 소리의 높이가 높다. 강세는 단어의 첫 음절에 위치하지만 모든 단어에 실현되지는 않고, 주로 명사와 동사 같은 내용어에 위치한다. 기본 높낮이는 2이고, 마지막 단어의 마지막 음절은 기본적으로 1이다. 평서문, 부정 명령문, 설명의문문에서는 강세 음절이 3을 받는다. 긍정 명령문에서는 내용어 외에 문장의 마지막 단어에도 강세가 오며 3을 받는다. 판정의문문에서는 강세 음절이 4를 받는다. 문장 마지막의 \term{e} `그렇지?'는 4를 받는다.

더 구체적으로, 동사와 명사의 첫 음절은 기본적으로 3이고 나머지 음절은 2이다. 대명사, 부사, 문법 기능어는 모든 음절이 기본적으로 2이다. 평서문 및 설명의문문에서는 마지막 단어가 동사나 명사라도 3이 오지 않으며 마지막 음절은 1로 하강한다. 명령문에서는 마지막 단어가 동사나 명사가 아니더라도 3이 오고 이후 하강한다. 판정의문문에서는 강세 음절이 3 대신 4를 받는다. \term{e} `그렇지?'가 붙은 의문문에서는 앞부분은 평서문/설명의문문과 같고 마지막 \term{e}가 4를 받는다.

\end{spacing}
\end{document}
